\section{Causal Inference}
We analyze causal models using potential outcomes, that is analyzing the causality between $Y_i$ and $D_i \in \{0, 1\}$. We call $Y_i(1)$ and $Y_i(0)$ to be the potential outcomes. \\
The \impt{treatment effect} is thus (causal effect of $D_i$): $Y_i(1) - Y_i(0)$. This leads to the \impt{average treatment effect} being
$$ATE = \expect{Y_i(1) - Y_i(0)}$$
We observe the following:
$$Y_i = Y_i(1)D_i + Y_i(0)(1 - D_i) = D_i(Y_i(1) - Y_i(0)) + Y_i(0)$$
Then, if we introduce $D_i$ to be treatment status, and $X_i$ to be some covariate, under the \impt{randomized control trial assumption}, that is $(Y_i(1), Y_i(0), X_i)$ are all independent of $D_i$, we have:
$$ATE = \frac{\Cov{Y_i}{D_i}}{\Var{D_i}} = \expect{Y_i \mid D_i = 1} - \expect{Y_i \mid D_i = 0} = \frac{\expect{Y_iD_i}}{\prob(D_i = 1)} - \frac{\expect{Y_i(1 - D_i)}}{\prob(D_i = 0)}$$
If we wish to test things about ATE, notice that we can take $W_i = \frac{Y_iD_i}{p} - \frac{Y_i(1 - D_i)}{1 - p}$, because we often know what $p$, and the test statistic would be
$$\tau = \frac{\sqrt{n}(\bar{W} - ATE_{H_0})}{\sqrt{\sVar{W_i}}}$$

\subsection{Selection on Observation}
Here, we no longer assume that the covariate $X_i$ is independent of $D_i$. However, we still need to make two assumptions to identify relevant parameters. \\
The first one is \impt{uncofoundedness} assumption, where $(Y_i(1), Y_i(0))$ are independent of $D_i$ given $X_i$. \\
The second one is \impt{overlap condition}, where
$$\exists \varepsilon > 0, \varepsilon < \prob(D_i = 1 \mid X_i = x) < 1 - \varepsilon$$
where we call $\prob(D_i = 1 \mid X_i = x)$ the \uimpt{propensity score}. \\
We introduce relevant properties of \impt{conditional independence}.
\begin{lemma}
    Suppose that $V$ and $W$ are independent given $X$, then for any function $\phi$:
    $$\expect{\phi(V) \mid W = w, X = x} = \expect{\phi(V) \mid X = x}$$
    If $Y$ is independent of $(X, W)$, then $Y$ is independent of $X$ given $W$. \\
    If $Y$ and $X$ are independent given $W$, then $f(Y)$ and $g(X)$ are independent given $W$.
\end{lemma}
The overlap condition implies that the $\text{supp}(X_i \mid D_i = 1) = \text{supp}(X_i \mid D_i = 0)$. \\
ATE is thus identified to be:
$$ATE = \frac{\expect{Y_iD_i}}{\prob(D_i = 1 \mid X_i)} - \frac{\expect{Y_i(1 - D_i)}}{\prob(D_i = 0 \mid X_i)}$$
This can take a sample analogue estimator given that defining $p(x) = \prob(D_i = 1 \mid X_i = x)$ and $\hat{p}(x) \pcvg p(x)$. This is an ``inverse propensity score weighted estimator''.

\subsection{Difference-in-differences}
Suppose we have $Y_{it}(d)$, where $t = 0, 1$ and $d = 0, 1$. We observe:
$$Y_{it} = D_i Y_{it}(1) + (1 - D_i)Y_{it}(0) = D_i(Y_{it}(1) - Y_{it}(0)) + Y_{it}(0)$$
and we define \impt{average treatment effect on the treated} to be:
$$ATT = \expect{Y_{i1}(1) - Y_{i1}(0) \mid D_i = 1}$$
$t = 0$ means before treatment is applied, and $t = 1$ means after treatment is applied. \\
We make two assumptions:
\begin{quote}
    1. No anticipation: $Y_{i0}(1) = Y_{i0}(0)$. \\
    2. Parallel trend assumption: 
    $$\expect{Y_{i1}(0) - Y_{i0}(0) \mid D_i = 1} = \expect{Y_{i1}(0) - Y_{i0}(0) \mid D_i = 0}$$
\end{quote}
We want to identify ATT. Note that by ``no anticipation'', we have $Y_{i0} = Y_{i0}(0)$. Then, by $Y_{i1} - Y_{i0}$ and take the expected value given $D_i = 0, 1$, we have:
$$ATT = \expect{Y_{i1} - Y_{i0} \mid D_i = 1} - \expect{Y_{i1} - Y_{i0} \mid D_i = 0}$$
Additionally, we have:
\begin{align*}
    ATT &= \expect{Y_{i1} - Y_{i0} \mid D_i = 1} - \expect{Y_{i1} - Y_{i0} \mid D_i = 0} \\
    &= \frac{\expect{\Delta Y_i D_i}}{\expect{D_i}} - \frac{\expect{\Delta Y_i (1 - D_i)}}{\expect{1 - D_i}}
\end{align*}
and then, we can provide a sample analogue estimator:
$$\hat{ATT} = \frac{\sCov{\Delta Y_i}{D_i}}{\sVar{D_i}}$$

\subsubsection{Two-way Fixed Effects}
Say we have:
$$Y_{it} = \alpha_i + \gamma_t + \beta \cdot 1\{t = 1\} \cdot D_i + u_{it}$$
with $\expect{u_{it} \mid \alpha_i, D_i} = 0$. $\beta$ can be identified in a similar way where:
$$\beta = \expect{\Delta Y_i \mid D_i = 1} - \expect{\Delta Y_i \mid D_i = 0}$$
As more multi-period two-way fixed effects, which is not discussed in this scope of the course, are discussed in many papers.

\subsubsection{Involvement of Covariate}
To do DID with covariance, we still have no anticipation assumption, but we update to \impt{conditional parallel trend assumption} where:
$$\expect{Y_{i1}(0) - Y_{i0}(0) \mid D_i = 1, X_i = x} = \expect{Y_{i1}(0) - Y_{i0}(0) \mid D_i = 0, X_i = x}$$
The conditional ATT is thus:
$$CATT = \expect{Y_{i1}(1) - Y_{i1}(0) \mid D_i = 1, X_i = x}$$
In this case, we call:
$$CATT(x) = \expect{\Delta Y_i \mid D_i = 1, X_i = x} - \expect{\Delta Y_i \mid D_i = 0, X_i = x}$$
Then, ATT is:
$$\expect{CATT(X_i) \mid D_i = 1} = \expect{\Delta Y_i \mid D_i = 1} - \expect{\expect{\Delta Y_i \mid D_i = 0, X_i} \mid D_i = 1}$$

\subsubsection{Estimating ATT: Linear Regression with Interaction Term}
Call
$$(Y_{10}, Y_{20}, \dots, Y_{n0}, Y_{11}, \dots, Y_{n1}) = (\tilde{Y}_1, \dots, \tilde{Y}_{2n})$$
Define $T_j$ to be $0$ if $\tilde{Y}_j$ comes from period $0$ and to be $1$ if $\tilde{Y}_j$ comes from period $1$. Define $\tilde{D}_j$ to be $0$ if $\tilde{Y}_j$ comes from the control group and to be $1$ otherwise. \\
We write $\gamma_{dt} = \expect{\tilde{Y}_j \mid \tilde{D}_j = d, T_j = t}$, and
$$\expect{\tilde{Y}_j \mid \tilde{D}_j, T_j} = \beta_0 + \beta_1 \tilde{D}_j + \beta_2 T_j  + \beta_3 \tilde{D}_j T_j$$
Then
$$\beta_3 = (\expect{\tilde{Y}_j \mid \tilde{D}_j = 1, T_j = 1} - \expect{\tilde{Y}_j \mid \tilde{D}_j = 1, T_j = 0}) - (\expect{\tilde{Y}_j \mid \tilde{D}_j = 0, T_j = 1} - \expect{\tilde{Y}_j \mid \tilde{D}_j = 0, T_j = 0})$$
This gives $\beta_3 = \gamma_{11} - \gamma_{10} - (\gamma_{01} - \gamma_{00})$. \\
Then $\beta_3 = ATT$ if we write $\tilde{Y}_j = \expect{\tilde{Y}_j \mid \tilde{D}_j, T_j} + u_j$


\newpage